% a06_BALtorus_ct1: Step 3(E), first part: the per-graph bound on the mass counterterm, |dm^2| <= c_m eps^{-1}
% (Ba2 p. 559, Ba3 (1.21)-(1.23)), by momentum-space power counting on the torus. Input from sec/a06_BALtorus_c.tex.

\paragraph{Step 3(E): the sizes of the mass counterterm.} Three statements about the mass counterterm $\dm$ are asserted in \cite{Ba2,Ba3} without proof and are consumed by every step of the chain: $\dm=O(\eps^{-1})$ \cite[p.~559]{Ba2}, and the two finite-part statements of \cite[p.~418]{Ba3}, $\dm_{K,k}(x)=O((L^k\eps)^{-1})$ and ``$\dm_{\mathrm{fin}}$ convergent as $\eps\to0$'' (``Future considerations will imply''), which enter the proof of Proposition~1 only through the size of the non-renormalising part of the \cite[(1.7)]{Ba3} vertex \cite[p.~421]{Ba3}. This step proves the first in the form stated, the second in the form stated, and the third in the form in which it is consumed: in scale-$k$ units the non-renormalising coefficient $c_k\defeq t_k^2(\dm_{\mathrm{fin}}+\dm_{K,k})$ is $O(t_k(1+\ln t_k^{-1})^5)$ with a constant free of $\eps_0$ (\cref{thm:a06e:T0}), which every consumer absorbs (\cref{lem:a06e:T1}); $\dm_{\mathrm{fin}}$ itself is bounded by a constant depending on $\eps_0$ (\cref{cor:a06e:T0p}).

\subsubsection*{The per-graph bound $\dm=O(\eps^{-1})$}

\emph{Setting} \cite[pp.~605--606]{Ba1}, \cite[pp.~413, 416--417]{Ba3}. $d=3$, $T_\eps\defeq\prod_\mu(\eps\Z/\ell_\mu\Z)$ with all sides $\ell_\mu\ge1$, $\ell_\mu/\eps\in\N$ (write $\ell_T^{-3}$ for $\prod_\mu\ell_\mu^{-1}$); functions carry the $\eps^3$-weighted product and $\hat f(p)\defeq\sum_x\eps^3f(x)e^{-ipx}$ for $p\in T^*\defeq\prod_\mu(2\pi/\ell_\mu)\Z\cap B_\eps$, $B_\eps\defeq[-\pi/\eps,\pi/\eps)^3$. For $p\in\R^3$, $[p]\in B_\eps$ is its representative modulo $(2\pi/\eps)\Z^3$. $(\partial_\mu f)(x)\defeq\eps^{-1}(f(x+\eps e_\mu)-f(x))$ has symbol $\hat\partial_\mu(p)=\eps^{-1}(e^{i\eps p_\mu}-1)$, and $\hat\Delta(p)\defeq\sum_\mu4\eps^{-2}\sin^2(\eps p_\mu/2)=\sum_\mu|\hat\partial_\mu(p)|^2$. The free covariances are $C_0\defeq(-\Delta+m^2)^{-1}$ (scalar, $m>0$) and $C\defeq(-\Delta+\lmu^2)^{-1}$ (vector, Feynman gauge, componentwise \cite[(1.11)]{Ba1}). By \cite[(1.20)--(1.21)]{Ba3}, $\dm=\sum_{2\le\alpha+2\beta\le4}e^\alpha\lambda^\beta\dm_{(\alpha,\beta)}$ is fixed graph by graph: $\dm_G\defeq\sum_{x'\in T_\eps}\eps^3\Sigma^\eps_G(x-x')$, $\Sigma^\eps_G$ the amputated one-particle-irreducible two-point graph $G$ (no difference on its external legs) with the vertices of \cite[(1.6)--(1.8), (1.10)]{Ba3} at $\eta=\eps$, $L^k\eps=1$, $\tilde B=0$, $g_k=1$ \cite[(1.22)--(1.23)]{Ba3}; $\dm$ is an $N\times N$ matrix polynomial in $q^2$, treated as a number. The vertices, the Taylor terms of $\frac12|D_A\phi|^2-\frac12|\partial\phi|^2=-\eps^{-2}\phi(b_-)\cdot(U(A_b)-1)\phi(b_+)$, are: $(V_\lambda)$ $-\lambda\sum_x\eps^3|\phi|^4$, order $2$; $(V_\delta)$ $-\frac12\sum_x\eps^3\dm_1|\phi|^2$ with $\dm_1$ the order-$2$ part of $\dm$, order $2$; $(V^8_n)$ $e^n(-1)^n\eps^{n-1}(n!)^{-1}\sum_b\eps^3[(\partial\phi)(b)\cdot q^n\phi(b_-)](A_b)^n$, $n\ge1$, order $n$; $(V^{10}_n)$ $e^n\eps^{n-2}(n!)^{-1}\sum_b\eps^3[\phi(b_-)\cdot q^n\phi(b_-)](A_b)^n$, $n\ge2$ even (for odd $n$, $q^n$ is antisymmetric and the term vanishes), order $n$. The order of $G$ is $\omega_G\defeq\sum_v\mathrm{ord}(v)=\alpha+2\beta$. The $\eps$-exponent of a vertex is $s_v\defeq0$ $(V_\lambda)$, $n-1$ $(V^8_n)$, $n-2$ $(V^{10}_n)$, and $s_v\defeq-1$ for $V_\delta$, the last used only after $|\dm_1|\le C\eps^{-1}$ is known; $s(G)\defeq\sum_vs_v$.

\begin{lemma}[symbols]\label{lem:a06e:M0}
$\frac4{\pi^2}|[p]|^2\le\hat\Delta(p)\le|[p]|^2$, $|\hat\partial_\mu(p)|\le|[p]|$, and $\hat C_0(p)=(\hat\Delta(p)+m^2)^{-1}\le\frac{\pi^2}4(|[p]|^2+m^2)^{-1}$, the same for $\hat C$ with $\lmu$.
\end{lemma}

\begin{proof}
All three are periodic, so take $p=[p]$; use $|\sin t|\ge2|t|/\pi$ for $|t|\le\pi/2$ and $|\sin t|\le|t|$; for the last claim, $\frac4{\pi^2}|p|^2+m^2\ge\frac4{\pi^2}(|p|^2+m^2)$.
\end{proof}

\begin{lemma}[Riemann domination]\label{lem:a06e:M1}
Let $F\colon B_\eps^{\,n}\to[0,\infty)$ satisfy $F(k)\le C_1F(k')$ whenever $|k_{i,\mu}-k'_{i,\mu}|\le\pi/\ell_\mu$ for all $i,\mu$. Then $C_1^{-1}(2\pi)^{-3n}\int_{B_\eps^n}F\le\ell_T^{-3n}\sum_{k\in(T^*)^n}F(k)\le C_1(2\pi)^{-3n}\int_{B_\eps^n}F$. For $F=\prod_cg_c(q_c(k))$ with $g_c(q)=(|[q]|^2+m_c^2)^{-\beta_c}$, $\beta_c\ge0$, each $q_c$ one of $\pm k_i$, $\pm(k_1+k_2)$, the hypothesis holds with $C_1=\prod_c(2+242/m_c^2)^{\beta_c}$.
\end{lemma}

\begin{proof}
The boxes of sides $2\pi/\ell_\mu$ centred at the points of $T^*$ tile $B_\eps$, and $\ell_T^{-3}=(2\pi)^{-3}|\text{box}|$; compare $F$ at the centre with $F$ on the cube. For the product form, $|[q_c(k)]-[q_c(k')]|\le2\sqrt3\pi/\min_\mu\ell_\mu\le11$, and for $|\delta|\le11$, $(|q'|+|\delta|)^2+m^2\le2|q'|^2+242+m^2\le(2+242/m^2)(|q'|^2+m^2)$.
\end{proof}

\begin{lemma}[radial integrals]\label{lem:a06e:M2}
For $\beta\in\R$, $m>0$, $0<\eps\le1$, $r>0$ let $I(\beta)\defeq\int_{B_\eps}(|p|^2+m^2)^{-\beta}dp$ and $H_\beta(r)\defeq\int_{B_\eps\cap\{|p|\le r\}}(|p|^2+m^2)^{-\beta}dp$. With $C=C(\beta,m)$: if $\beta<\frac32$, $I(\beta)\le C\eps^{2\beta-3}$ and $H_\beta(r)\le C(1+r)^{3-2\beta}$; if $\beta=\frac32$, $I\le C(1+\ln\eps^{-1})$ and $H\le C(1+\ln(1+r))$; if $\beta>\frac32$, $I,H\le C$.
\end{lemma}

\begin{proof}
$B_\eps\subset\{|p|\le\sqrt3\pi/\eps\}$; in polar coordinates $(s^2+m^2)^{-\beta}\le\max(m^{-2\beta},(1+m^2)^{-\beta})$ on $s\le1$ and $\le Cs^{-2\beta}$ on $s\ge1$; $\int_1^{\max(\rho,1)}s^{2-2\beta}ds$ gives the three cases.
\end{proof}

\begin{lemma}[structure and degree]\label{lem:a06e:M3}
Let $G$ be a graph of $\dm$ of order $\omega_G\le4$. Then: \textup{(a)} the loop number is $\mathcal L\defeq I-V+1=(\omega_G-\omega_{\mathrm{ins}})/2$, $\omega_{\mathrm{ins}}$ the total order of the $V_\delta$ vertices; so $\mathcal L\le2$, and a $V_\delta$ vertex occurs only in one-loop graphs of order $4$, exactly one, of order $2$. \textup{(b)} The internal-line graph $G'$ is connected and bridgeless; for $\mathcal L=1$ it is a cycle, for $\mathcal L=2$ a theta (two vertices of degree $3$ joined by three internally disjoint chains) or a figure-eight (one vertex of degree $4$ on two cycles); all other vertices have degree $2$ in $G'$. \textup{(c)} At zero external momentum all lines of a chain (or cycle) carry one momentum $q_c$ ($k_1,k_2,k_1+k_2$ for a theta; $k_1,k_2$ for a figure-eight; $k_1$ for a cycle), and $|\dm_G|\le C_G\eps^{s(G)}\ell_T^{-3\mathcal L}\sum_{k\in(T^*)^{\mathcal L}}\prod_cg_c(q_c)$ with $g_c(q)=(|[q]|^2+m_c^2)^{-\beta_c}$, $\beta_c\defeq\sum_{l\in c}(1-d_l/2)\ge0$, $d_l\in\{0,1,2\}$ the number of differences acting on the line $l$, $m_c\in\{m,\lmu\}$ on a chain of one line type and $m_c\defeq\min(m,\lmu)$ on a mixed chain. \textup{(d)} With $\omega(G)\defeq3\mathcal L-2\sum_l(1-d_l/2)-s(G)$ and, for a cycle $H$ of $G'$, $\omega(H)\defeq3-2\sum_{l\in H}(1-d_l/2)-s(H)$: $\omega(G)=2-\omega_G/2$, and every cycle $H\ne G'$ of a two-loop $G'$ has $\omega(H)\le1$, with equality only if $H$ is a two-point subgraph of order $2$.
\end{lemma}

\begin{proof}
(a) $\sum_vh_v=2I+2$ ($h_v$ the number of legs), so $\mathcal L=\sum_v(h_v/2-1)$, and $h_v/2-1=\mathrm{ord}(v)/2$ for $V_\lambda$ ($h=4$), $V^8_n$, $V^{10}_n$ ($h=n+2$) and $0$ for $V_\delta$ ($h=2$). (b) $G$ connected with external legs at vertices gives $G'$ connected; 1PI gives no bridge, so no vertex of degree $\le1$; $\sum_v(\deg_{G'}v-2)=2(\mathcal L-1)$. For $\mathcal L=2$ with two vertices $u,w$ of degree $3$, a path leaving $u$ and returning to $u$ would make the connection to $w$ a bridge, so all three paths end at $w$. (c) Momentum conservation at the vertices with zero external momenta; per line $|\hat C(q)||\hat\partial(q)|^{d_l}\le\frac{\pi^2}4(|[q]|^2+m_c^2)^{-1+d_l/2}$ by \cref{lem:a06e:M0}; vertex phases have modulus $1$, $\norm q\le1$. (d) With $\mathcal L=I-V+1$ and $2I=\sum h_v-2$, $\omega(G)=\sum_v(h_v/2-3-s_v+d_v)+2$, $d_v=1$ for $V^8$ vertices (the difference acts on an internal line) and $0$ otherwise; the bracket is $-\mathrm{ord}(v)/2$ for every type ($-1$, $-1$, $(n+2)/2-3-(n-1)+1=-n/2$, $(n+2)/2-3-(n-2)=-n/2$). For a cycle $H$ (as many lines as vertices), $\omega(H)=3-\mathrm{ord}(H)/2-E_H/2-\#\{V^8\text{ vertices of }H\text{ whose difference acts off }H\}$, $E_H\defeq\sum_{v\in H}(h_v-2)$: the bracket per vertex is $-1-e_v/2$ ($V_\lambda$, $e_v=2$), $-n/2-e_v/2-(1-d_{v,H})$ ($V^8_n$), $-n/2-e_v/2$ ($V^{10}_n$), $-1-e_v/2$ ($V_\delta$). For a proper cycle of a theta or a figure-eight, $E_H\ge2$ (the third chain, resp.\ the other cycle, leaves $H$). $\mathrm{ord}(H)\ge2$: a cycle through $\ge2$ vertices has order $\ge2$; a self-loop can sit only at the junction of a figure-eight, which has $\ge4$ internal legs, so it is $V_\lambda$, $V^8_n$ or $V^{10}_n$ with $n\ge2$. Hence $\omega(H)\le1$, with equality iff $\mathrm{ord}(H)=2$, $E_H=2$ and no difference leaves $H$.
\end{proof}

\begin{lemma}[one loop]\label{lem:a06e:M4}
If $\mathcal L(G)=1$, then $|\dm_G|\le C\eps^{s(G)}I(\beta)$ with $\beta$ the $\beta_c$ of the cycle, and $|\dm_G|\le C\eps^{-1}$.
\end{lemma}

\begin{proof}
The first bound is \cref{lem:a06e:M1,lem:a06e:M3}(c). Then $\omega(G)=3-2\beta-s$ and \cref{lem:a06e:M2}. Order $2$ (no $V_\delta$, $s\ge0$, $\omega=1$): $s=0$ gives $\beta=1$, $I\le C\eps^{-1}$; $s\ge1$ gives $\beta=1-s/2<\frac32$ and $\eps^sI\le C\eps^{s+2\beta-3}=C\eps^{-1}$. Order $4$ with one $V_\delta$ ($s=s'-1$, $s'\ge0$ the sum over the other vertices, $\omega=0$, $3-2\beta=s'-1$): $s'=0$ gives $\beta=2$ and $C\eps^{-1}$; $s'=1$ gives $\beta=\frac32$ and $C(1+\ln\eps^{-1})$; $s'\ge2$ gives $\eps^{s'-1}\eps^{2\beta-3}=1$.
\end{proof}

\begin{lemma}[theta]\label{lem:a06e:M5}
If $G'$ is a theta with chains $c=1,2,3$, then $|\dm_G|\le C\eps^{-1}$.
\end{lemma}

\begin{proof}
By \cref{lem:a06e:M1,lem:a06e:M3}(c), $|\dm_G|\le C\eps^s\,T$ with $T\defeq\iint_{B_\eps^2}g_1(k)g_2(l)g_3(k+l)\,dk\,dl$; at order $4$ with $\mathcal L=2$ there is no $V_\delta$, so $s\ge0$. (1) For $a\in\{1,2,3\}$ let $S_a$ be the set where $|[q_a]|$ is the smallest of $|[k]|,|[l]|,|[k+l]|$, $b,c$ the other indices and $q_*$ the one of $q_b,q_c$ of smaller norm; each $g$ decreases in the norm, so on $S_a$ the integrand is $\le g_a(q_a)g_b(|q_*|)g_c(|q_*|)$ with $|[q_a]|\le|[q_*]|$. (2) Any two of $k,l,k+l$ are coordinates with unit Jacobian, so $\iint_{S_a}\le2\int_{B_\eps}g_b(q)g_c(q)H_{\beta_a}(|q|)dq$. (3) Put $\omega^0\defeq6-2(\beta_1+\beta_2+\beta_3)$, $\omega^0_{bc}\defeq3-2(\beta_b+\beta_c)$; then $\omega(G)=\omega^0-s=0$ and $\omega(H_{bc})=\omega^0_{bc}-s(H_{bc})\le1$ by \cref{lem:a06e:M3}(d), with $s(H_{bc})\le s$. (4) $\beta_a<\frac32$: $H\le C(1+|q|)^{3-2\beta_a}\le C(|q|^2+m^2)^{(3-2\beta_a)/2}$, so the sector is $\le C\,I(\beta')$ with $3-2\beta'=\omega^0=s$, and $\eps^s$ times it is $\le C(1+\ln\eps^{-1})$. (5) $\beta_a=\frac32$: $H\le C(1+\ln\eps^{-1})$ and $\omega^0_{bc}=\omega^0=s$, so the sector is $\le C(1+\ln\eps^{-1})^2$. (6) $\beta_a>\frac32$: $H\le C$ and $\eps^sI(\beta_b+\beta_c)\le C\max(\eps^{s-\omega^0_{bc}},1+\ln\eps^{-1})$ with $\eps^{s-\omega^0_{bc}}\le\eps^{s(H_{bc})-\omega^0_{bc}}=\eps^{-\omega(H_{bc})}\le\eps^{-1}$. Sum the six pieces.
\end{proof}

\begin{lemma}[figure-eight]\label{lem:a06e:M6}
If $G'$ is a figure-eight with cycles $1,2$, then $|\dm_G|\le C\eps^sI(\beta_1)I(\beta_2)\le C\eps^{-1}$.
\end{lemma}

\begin{proof}
The cycles carry independent momenta; $\omega(G)=6-2\beta_1-2\beta_2-s=0$ and $\omega(H_i)\le1$. Both $\beta_i<\frac32$: $\eps^{s+2\beta_1-3+2\beta_2-3}=\eps^{-\omega(G)}=1$. $\beta_1<\frac32=\beta_2$: $\eps^{-\omega(G)}(1+\ln\eps^{-1})$. $\beta_1<\frac32<\beta_2$: $\eps^{s-(3-2\beta_1)}\le\eps^{-\omega(H_1)}\le\eps^{-1}$, using $s\ge s(H_1)$. Both $\ge\frac32$: $\le C(1+\ln\eps^{-1})^2$.
\end{proof}

\begin{theorem}[the mass counterterm per graph]\label{thm:a06e:CT1}
For every graph $G$ of $\dm$ (\cite[(1.22)--(1.23)]{Ba3}: all graphs of order $\alpha+2\beta\le4$, $d=3$), $|\dm_G|\le C_G\eps^{-1}$ with $C_G=C_G(m,\lmu,N)$, uniformly in $0<\eps\le1$ and in the torus (all sides $\ge1$). Hence $|\dm|\le c_m\eps^{-1}$ with $c_m\defeq\sum_Ge^\alpha\lambda^\beta C_G$, a finite sum; $c_m$ does not depend on $\eps_0$.
\end{theorem}

\begin{proof}
Order $2$: by \cref{lem:a06e:M3}(a) these are one-loop graphs without $V_\delta$, so \cref{lem:a06e:M4} gives $|\dm_1|\le C\eps^{-1}$, which justifies $s_v=-1$ for the $V_\delta$ vertices at order $4$. Order $4$: by \cref{lem:a06e:M3} each graph is a cycle with one $V_\delta$ (\cref{lem:a06e:M4}), a theta (\cref{lem:a06e:M5}) or a figure-eight (\cref{lem:a06e:M6}). At order $\le4$, $n\le4$ and $V\le4$, so there are finitely many graphs.
\end{proof}

The proof uses the momentum-space form: in $x$-space the photon--scalar bubble line $\partial C_0\partial^*$ has a logarithmically divergent $\ell^1$ norm, while its symbol $|\hat\partial|^2\hat C_0$ is bounded (\cref{lem:a06e:M0}). Not every order-$4$ graph is $O(\ln\eps^{-1})$: graphs with an order-$2$ two-point subgraph (the last graph of \cite[(1.22)]{Ba3}, the tadpole on a tadpole of \cite[(1.23)]{Ba3}) are of size $\eps^{-1}$. The third graph of \cite[(1.22)]{Ba3} is \cref{lem:a06e:M6} at the junction $V^{10}_4$: $\eps^2C_0(0)^2=O(1)$, Ba{\l}aban's ``except the third''.

\begin{corollary}[{\cite[(2.11)--(2.13)]{Ba2}}]\label{cor:a06e:C1}
The bounds \cite[(2.11)--(2.13)]{Ba2} on the counterterm vertices hold for Ba{\l}aban's torus $\dm$ with constants free of $\eps$, $\eps_0$ and the torus.
\end{corollary}

\begin{proof}
They are stated in Step~3 as consequences of $|\dm|\le c\eps^{-1}$ by exact computation; \cref{thm:a06e:CT1} supplies the hypothesis.
\end{proof}
